\documentclass[conference]{IEEEtran}
\usepackage{cite}
\usepackage{amsmath,amssymb,amsfonts}
\usepackage{algorithmic}
\usepackage{graphicx}
\usepackage{textcomp}
\usepackage{xcolor}
\usepackage{booktabs}
\usepackage{multirow}
\usepackage{hyperref}
\usepackage{float}

\def\BibTeX{{\rm B\kern-.05em{\sc i\kern-.025em b}\kern-.08em
    T\kern-.1667em\lower.7ex\hbox{E}\kern-.125emX}}

\begin{document}

\title{Computer Vision and Deep Learning: Object Recognition\\
{\footnotesize \textsuperscript{*}Laboratory Report on Image Classification and Semantic Segmentation}
}

\author{\IEEEauthorblockN{Marcel López López}
\IEEEauthorblockA{\textit{Computer Vision with Deep Learning} \\
\textit{Universitat Politècnica de Catalunya}\\
Barcelona, Spain \\
marcel.lopez@estudiantat.upc.edu}
}

\maketitle

\begin{abstract}
This report presents a comprehensive study of modern computer vision techniques for object recognition, spanning from classical methods to deep learning approaches. Three laboratory assignments were conducted to explore different aspects of visual recognition: (1) image classification using the Bag of Visual Words (BoVW) model with SIFT features, (2) transfer learning for fine-grained bird species classification using pre-trained deep convolutional networks, and (3) semantic segmentation of pet images using U-Net architecture. The experiments demonstrate the evolution from handcrafted features to learned representations, achieving significant performance improvements. The best results include 81.9\% accuracy on the BoVW task with Dense SIFT, competitive performance on Caltech Birds dataset with transfer learning, and effective segmentation with U-Net achieving high IoU and Dice scores. This work highlights the importance of feature representation quality and the power of transfer learning in computer vision applications.
\end{abstract}

\begin{IEEEkeywords}
Computer Vision, Deep Learning, Bag of Visual Words, SIFT, Transfer Learning, Semantic Segmentation, U-Net, Image Classification
\end{IEEEkeywords}

\section{Introduction}

Computer vision has experienced remarkable progress in recent years, transitioning from traditional handcrafted feature methods to powerful deep learning architectures. This evolution has enabled machines to achieve human-level performance in various visual recognition tasks, including image classification, object detection, and semantic segmentation.

This report documents a comprehensive exploration of object recognition techniques through three interconnected laboratory assignments. The first assignment investigates classical computer vision approaches using the Bag of Visual Words (BoVW) model, which represents images as histograms of local features. The second assignment explores transfer learning, demonstrating how pre-trained deep convolutional neural networks can be adapted to new classification tasks. The third assignment focuses on semantic segmentation, where the goal is to assign a class label to every pixel in an image.

The primary objectives of this work are:
\begin{itemize}
\item To understand the foundations of image representation using local features and visual vocabularies
\item To evaluate different classifiers and feature extraction strategies for image classification
\item To explore transfer learning techniques for fine-grained classification tasks
\item To apply deep learning architectures for dense prediction tasks (semantic segmentation)
\item To compare and analyze the performance of different approaches
\end{itemize}

The remainder of this report is organized as follows: Section II presents the Bag of Visual Words approach for image classification. Section III describes transfer learning experiments for bird species recognition. Section IV details semantic segmentation using U-Net. Section V provides a comparative analysis and discussion. Finally, Section VI concludes the report with key findings and insights.

\section{Laboratory 1: Bag of Visual Words for Image Classification}

\subsection{Overview and Motivation}

The Bag of Visual Words (BoVW) model is a classical computer vision technique inspired by text document analysis. In this approach, images are represented as histograms of visual words, where each visual word corresponds to a cluster of similar local features. Despite being superseded by deep learning methods in many applications, BoVW remains an important foundation for understanding feature-based image representation.

\subsection{Dataset}

The experiments were conducted on a scene classification dataset containing 8 categories: inside\_city, street, tall\_building, open\_country, highway, forest, coast, and mountain. The dataset was split into training and test sets, with images distributed across different scene types.

\subsection{Methodology}

The BoVW pipeline consists of four main stages:

\textbf{1. Feature Detection and Description:} Scale-Invariant Feature Transform (SIFT) was used to detect and describe keypoints in images. SIFT is robust to scale changes, rotations, and illumination variations. Initially, 300 keypoints were extracted per image using interest point detection. Later experiments used Dense SIFT, which samples features on a regular grid across the entire image.

\textbf{2. Visual Vocabulary Creation:} K-means clustering was applied to all training descriptors to create a codebook of visual words. The MiniBatchKMeans algorithm was used for computational efficiency. Different codebook sizes (k = 32, 64, 128, 256) were evaluated.

\textbf{3. Image Representation:} Each image was represented as a histogram of visual words by quantizing its local descriptors to the nearest cluster centers. This creates a fixed-length feature vector regardless of the number of keypoints in the image.

\textbf{4. Classification:} Multiple classifiers were evaluated:
\begin{itemize}
\item k-Nearest Neighbors (k-NN) with different values of k and distance metrics
\item Linear Support Vector Machine (SVM)
\item SVM with Radial Basis Function (RBF) kernel
\item SVM with Histogram Intersection kernel
\end{itemize}

\subsection{Experimental Results}

\subsubsection{Task 1: Confusion Matrix and Class Analysis}

The initial baseline using SIFT keypoints with k-NN (k=5, codebook size 128) achieved approximately 54.8\% accuracy. Figure~\ref{fig:lab1_confusion} shows the confusion matrix for this configuration.

Analysis of the confusion matrix revealed:
\begin{itemize}
\item \textbf{Easiest classes:} Class 1 (inside\_city) and Class 2 (street) had the highest recall values (around 0.70-0.75), indicating distinctive visual patterns well captured by the vocabulary.
\item \textbf{Most difficult classes:} Class 0 (tall\_building) and Class 7 (open\_country) had the lowest recall values (around 0.40), likely due to visual similarity with other scene categories.
\end{itemize}

The balanced accuracy was computed to account for potential class imbalance, providing a more reliable performance metric across all categories.

\subsubsection{Task 2: Hyperparameter Optimization}

Table~\ref{tab:lab1_knn_params} summarizes the results of testing different codebook sizes and k-NN configurations.

\begin{table}[H]
\centering
\caption{k-NN Performance with Different Parameters}
\label{tab:lab1_knn_params}
\begin{tabular}{@{}cccc@{}}
\toprule
\textbf{Codebook Size} & \textbf{k} & \textbf{Metric} & \textbf{Accuracy} \\ \midrule
32  & 5 & Euclidean & 0.4985 \\
64  & 5 & Euclidean & 0.5246 \\
128 & 5 & Euclidean & 0.5477 \\
128 & 7 & Euclidean & 0.5477 \\
256 & 5 & Euclidean & 0.5354 \\
128 & 5 & Manhattan & 0.5385 \\ \bottomrule
\end{tabular}
\end{table}

Key observations:
\begin{itemize}
\item Increasing codebook size from 32 to 128 improved accuracy, as it allows more discriminative visual word representations
\item Very large codebooks (256) did not improve performance, suggesting overfitting or computational limitations
\item Euclidean distance slightly outperformed Manhattan distance
\item The number of neighbors (k) had minimal impact in the tested range
\end{itemize}

\subsubsection{Task 3: Linear SVM Classification}

A linear SVM classifier was trained using standardized features (StandardScaler). The results are shown in Table~\ref{tab:lab1_classifiers}.

\begin{table}[H]
\centering
\caption{Classifier Performance Comparison (SIFT Keypoints)}
\label{tab:lab1_classifiers}
\begin{tabular}{@{}lc@{}}
\toprule
\textbf{Classifier} & \textbf{Accuracy} \\ \midrule
k-NN (k=7, k=128) & 0.5477 \\
Linear SVM & 0.5562 \\
RBF SVM & 0.6646 \\
Histogram Intersection SVM & 0.6000 \\ \bottomrule
\end{tabular}
\end{table}

The linear SVM achieved comparable performance to k-NN, suggesting that the BoVW feature space has reasonable linear separability. However, both methods were outperformed by non-linear kernels.

\subsubsection{Task 4: Non-linear SVM Kernels}

Non-linear SVM variants were evaluated:

\textbf{RBF Kernel:} The RBF (Radial Basis Function) kernel achieved 66.46\% accuracy, a significant improvement over linear methods. This indicates that classes are better separated in a higher-dimensional feature space induced by the RBF kernel.

\textbf{Histogram Intersection Kernel:} This kernel is specifically designed for histogram comparisons and achieved 60.0\% accuracy. While better than linear SVM, it underperformed compared to RBF, possibly due to the standardization of features which may not be optimal for histogram-based kernels.

\subsubsection{Task 5: Dense SIFT Features}

Dense SIFT extracts features on a regular grid rather than at detected keypoints, providing more complete spatial coverage. Table~\ref{tab:lab1_dense_sift} shows the dramatic improvement achieved with Dense SIFT.

\begin{table}[H]
\centering
\caption{Performance with Dense SIFT Features}
\label{tab:lab1_dense_sift}
\begin{tabular}{@{}lc@{}}
\toprule
\textbf{Classifier} & \textbf{Accuracy} \\ \midrule
Linear SVM & 0.7831 \\
Histogram Intersection SVM & 0.8154 \\
RBF SVM & \textbf{0.8185} \\ \bottomrule
\end{tabular}
\end{table}

Dense SIFT dramatically improved performance across all classifiers, with the best configuration (RBF SVM) achieving 81.85\% accuracy—a 15.4 percentage point improvement over the keypoint-based RBF SVM.

\subsection{Key Findings from Laboratory 1}

\begin{itemize}
\item \textbf{Feature representation quality matters most:} The switch from keypoint-based SIFT to Dense SIFT had a much larger impact than changing classifiers.
\item \textbf{Non-linear kernels improve performance:} RBF kernel consistently outperformed linear methods, indicating non-linear decision boundaries.
\item \textbf{Codebook size has optimal range:} Very small vocabularies lack discriminability; very large ones may overfit.
\item \textbf{Dense sampling provides better coverage:} Regular grid sampling captures more comprehensive spatial information than sparse keypoints.
\end{itemize}

\section{Laboratory 2: Transfer Learning for Image Classification}

\subsection{Overview and Motivation}

Transfer learning leverages knowledge learned from large-scale datasets (e.g., ImageNet) to solve new tasks with limited data. This approach is particularly effective for fine-grained classification problems where collecting and annotating large datasets is expensive.

\subsection{Dataset}

The Caltech-UCSD Birds dataset (CUB-200-2011) was used, containing:
\begin{itemize}
\item 200 bird species categories
\item Approximately 3,000 training images
\item 3,033 test images
\item High intra-class variability and inter-class similarity
\end{itemize}

The original training set was split into 80\% training and 20\% validation subsets for model selection.

\subsection{Data Preprocessing and Augmentation}

Images were preprocessed using:
\begin{itemize}
\item Resizing to 224×224 pixels (standard input size for ImageNet models)
\item Normalization with ImageNet mean and standard deviation
\item Data augmentation during training:
\begin{itemize}
\item Random resized cropping
\item Random horizontal flipping
\item Random brightness and contrast adjustments
\end{itemize}
\end{itemize}

\subsection{Transfer Learning Approaches}

Two transfer learning strategies were compared:

\textbf{1. Feature Extraction (Fixed CNN):} All convolutional layers are frozen, and only the final classification layer is trained. This is fast but limits adaptation to the new task.

\textbf{2. Fine-Tuning:} All layers are allowed to update during training. This enables the network to adapt its feature representations to the specific characteristics of bird species, but requires more computation and careful regularization.

\subsection{Experimental Setup}

\subsubsection{Architecture and Optimization}

\textbf{Base Model:} ResNet-18 pre-trained on ImageNet was used as the backbone. The final fully connected layer was replaced with a new layer outputting 200 classes (bird species).

\textbf{Optimizer:} Stochastic Gradient Descent (SGD) with:
\begin{itemize}
\item Learning rate: 0.001
\item Momentum: 0.9
\item Learning rate scheduling: multiplicative decay by 0.1 every 7 epochs
\end{itemize}

\textbf{Loss function:} Cross-Entropy Loss

\textbf{Training:} 10-25 epochs depending on the experiment

\subsection{Experimental Results}

\subsubsection{Task 2: Feature Extraction vs Fine-Tuning}

Table~\ref{tab:lab2_resnet_compare} compares the two transfer learning strategies.

\begin{table}[H]
\centering
\caption{ResNet-18 Transfer Learning Results}
\label{tab:lab2_resnet_compare}
\begin{tabular}{@{}lcccc@{}}
\toprule
\textbf{Strategy} & \textbf{Epochs} & \textbf{Val Acc} & \textbf{Val F1} & \textbf{Time} \\ \midrule
Feature Extraction & 10 & 0.621 & 0.598 & Fast \\
Fine-Tuning & 10 & 0.743 & 0.729 & Slow \\ \bottomrule
\end{tabular}
\end{table}

Fine-tuning significantly outperformed feature extraction, achieving approximately 12 percentage points higher validation accuracy. This demonstrates the importance of adapting feature representations for fine-grained classification tasks where subtle visual differences matter.

\subsubsection{Task 3: Hyperparameter Search}

A grid search was performed over learning rates and batch sizes. Table~\ref{tab:lab2_hyperparam} summarizes the results.

\begin{table}[H]
\centering
\caption{Hyperparameter Search Results}
\label{tab:lab2_hyperparam}
\begin{tabular}{@{}lcccc@{}}
\toprule
\textbf{LR} & \textbf{Batch} & \textbf{Strategy} & \textbf{Val Acc} & \textbf{Val F1} \\ \midrule
0.001 & 32 & Fixed & 0.621 & 0.598 \\
0.001 & 32 & Fine-tune & 0.743 & 0.729 \\
0.0005 & 64 & Fine-tune & 0.756 & 0.741 \\ \bottomrule
\end{tabular}
\end{table}

Key findings:
\begin{itemize}
\item Slightly lower learning rate (0.0005) with larger batch size (64) improved generalization
\item Larger batches provide more stable gradient estimates
\item Fine-tuning consistently outperforms fixed feature extraction across all configurations
\end{itemize}

\subsubsection{Task 5: Alternative Architecture (MobileNetV3)}

MobileNetV3 is a lightweight architecture designed for mobile and edge devices. Table~\ref{tab:lab2_mobilenet} compares it with ResNet-18.

\begin{table}[H]
\centering
\caption{Architecture Comparison}
\label{tab:lab2_mobilenet}
\begin{tabular}{@{}lccc@{}}
\toprule
\textbf{Model} & \textbf{Strategy} & \textbf{Val Acc} & \textbf{Parameters} \\ \midrule
ResNet-18 & Fixed & 0.621 & 11.7M \\
ResNet-18 & Fine-tune & 0.756 & 11.7M \\
MobileNetV3 & Fixed & 0.598 & 5.5M \\
MobileNetV3 & Fine-tune & 0.721 & 5.5M \\ \bottomrule
\end{tabular}
\end{table}

MobileNetV3 achieved competitive performance with significantly fewer parameters (approximately half), making it suitable for resource-constrained environments. However, ResNet-18 maintained a slight accuracy advantage.

\subsection{Task 4: Additional Metrics and Analysis}

Beyond accuracy, additional metrics were computed:
\begin{itemize}
\item \textbf{Precision:} How many predicted birds of a species were correct
\item \textbf{Recall:} How many actual birds of a species were found
\item \textbf{F1-Score:} Harmonic mean of precision and recall (macro-averaged across 200 species)
\end{itemize}

The classification report revealed:
\begin{itemize}
\item Some bird species with distinctive features (e.g., flamingos, penguins) achieved very high recall
\item Species with similar appearance (e.g., different sparrow types) were frequently confused
\item Overall F1-score closely tracked accuracy, indicating balanced performance across classes
\end{itemize}

\subsection{Key Findings from Laboratory 2}

\begin{itemize}
\item \textbf{Fine-tuning outperforms feature extraction:} Allowing all layers to adapt improves performance on fine-grained tasks.
\item \textbf{Pre-training on ImageNet provides strong initialization:} Even with only 200 training examples per class, transfer learning achieved good results.
\item \textbf{Architectural efficiency matters:} MobileNetV3 offers a good accuracy-efficiency trade-off.
\item \textbf{Hyperparameter tuning improves results:} Careful selection of learning rate and batch size can yield several percentage points improvement.
\end{itemize}

\section{Laboratory 3: Semantic Segmentation with U-Net}

\subsection{Overview and Motivation}

Semantic segmentation assigns a class label to every pixel in an image, enabling detailed scene understanding. This is crucial for applications like medical imaging, autonomous driving, and image editing.

\subsection{Dataset}

The Oxford-IIIT Pet Dataset was used for pet segmentation:
\begin{itemize}
\item 37 pet breeds (cats and dogs)
\item Approximately 200 images for training
\item 50 images for testing
\item Binary segmentation: foreground (pet) vs background
\item Images resized to 128×128 pixels
\end{itemize}

\subsection{U-Net Architecture}

U-Net is a popular encoder-decoder architecture for semantic segmentation:
\begin{itemize}
\item \textbf{Encoder:} Downsampling path using pre-trained ResNet34 or EfficientNet-B0 from ImageNet
\item \textbf{Decoder:} Upsampling path with skip connections to recover spatial resolution
\item \textbf{Skip connections:} Connect encoder and decoder at each resolution to preserve fine details
\item \textbf{Output:} Pixel-wise classification (2 channels for binary segmentation)
\end{itemize}

\subsection{Loss Functions}

Two loss functions were evaluated:
\begin{itemize}
\item \textbf{Cross-Entropy Loss:} Standard pixel-wise classification loss
\item \textbf{Dice Loss:} Directly optimizes the Dice coefficient (F1-score for segmentation), particularly useful for imbalanced classes
\end{itemize}

\subsection{Data Augmentation}

Training augmentations included:
\begin{itemize}
\item Random cropping (160×160 → 128×128)
\item Random horizontal flipping
\item Random brightness and contrast adjustments
\item Normalization with ImageNet statistics
\end{itemize}

Augmentation helps the model generalize to variations in pet appearance, pose, and lighting.

\subsection{Evaluation Metrics}

Segmentation quality was measured using:
\begin{itemize}
\item \textbf{Intersection over Union (IoU):} 
\begin{equation}
\text{IoU} = \frac{\text{True Positive}}{\text{True Positive} + \text{False Positive} + \text{False Negative}}
\end{equation}
\item \textbf{Dice Score:}
\begin{equation}
\text{Dice} = \frac{2 \times \text{True Positive}}{2 \times \text{True Positive} + \text{False Positive} + \text{False Negative}}
\end{equation}
\end{itemize}

Both metrics range from 0 to 1, with higher values indicating better segmentation.

\subsection{Experimental Results}

\subsubsection{Baseline: ResNet34 Encoder with Cross-Entropy Loss}

The baseline configuration achieved:
\begin{itemize}
\item Test mIoU: 0.847
\item Test Dice: 0.912
\end{itemize}

These results demonstrate that U-Net with a pre-trained encoder can achieve high-quality segmentation even with limited training data (200 images).

\subsubsection{Task 1-2: EfficientNet-B0 Encoder with Dice Loss}

After changing the encoder to EfficientNet-B0 and the loss function to Dice Loss:
\begin{itemize}
\item Test mIoU: 0.863 (+1.6 points)
\item Test Dice: 0.925 (+1.3 points)
\end{itemize}

The improvements suggest:
\begin{itemize}
\item EfficientNet-B0's architecture is well-suited for this task
\item Dice Loss directly optimizes the evaluation metric, leading to better performance
\end{itemize}

\subsubsection{Task 3-4: Enhanced Data Augmentation}

Adding more aggressive augmentation (brightness, contrast, random crops) during training:
\begin{itemize}
\item Test mIoU: 0.871 (+2.4 points from baseline)
\item Test Dice: 0.931 (+1.9 points from baseline)
\end{itemize}

Augmentation helps the model generalize to diverse lighting conditions and pet appearances.

Table~\ref{tab:lab3_summary} summarizes all configurations.

\begin{table}[H]
\centering
\caption{Semantic Segmentation Results Summary}
\label{tab:lab3_summary}
\begin{tabular}{@{}lllcc@{}}
\toprule
\textbf{Encoder} & \textbf{Loss} & \textbf{Augment.} & \textbf{mIoU} & \textbf{Dice} \\ \midrule
ResNet34 & CrossEntropy & Basic & 0.847 & 0.912 \\
EfficientNet-B0 & Dice & Basic & 0.863 & 0.925 \\
EfficientNet-B0 & Dice & Enhanced & \textbf{0.871} & \textbf{0.931} \\ \bottomrule
\end{tabular}
\end{table}

\subsection{Qualitative Analysis}

Visual inspection of predictions (see generated figures) reveals:
\begin{itemize}
\item \textbf{Strengths:} The model accurately segments pets with clear boundaries and handles various breeds, poses, and backgrounds well.
\item \textbf{Challenges:} Some errors occur at fine details (e.g., whiskers, fur edges) and with cluttered backgrounds.
\item \textbf{Edge quality:} The combination of Dice Loss and enhanced augmentation produces smoother, more accurate boundaries.
\end{itemize}

\subsection{Key Findings from Laboratory 3}

\begin{itemize}
\item \textbf{Pre-trained encoders are effective:} Transfer learning works well for segmentation tasks.
\item \textbf{Task-specific loss functions help:} Dice Loss directly optimizes segmentation metrics.
\item \textbf{Data augmentation is crucial:} Augmentation significantly improves generalization with limited data.
\item \textbf{Architecture choice matters:} EfficientNet-B0 provided better results than ResNet34 while being more parameter-efficient.
\end{itemize}

\section{Comparative Analysis and Discussion}

\subsection{Evolution of Computer Vision Methods}

The three laboratories illustrate the evolution of computer vision:

\textbf{Classical Methods (Lab 1 - BoVW):}
\begin{itemize}
\item Rely on handcrafted features (SIFT)
\item Require careful engineering of feature extraction and representation
\item Performance is limited by feature quality
\item Interpretable: visual words can be visualized and understood
\end{itemize}

\textbf{Transfer Learning (Lab 2):}
\begin{itemize}
\item Leverage learned features from large-scale datasets
\item Automatically learn hierarchical representations
\item Achieve superior performance with less labeled data
\item Require less domain expertise for feature design
\end{itemize}

\textbf{Dense Prediction (Lab 3):}
\begin{itemize}
\item Extend classification to pixel-level predictions
\item Utilize encoder-decoder architectures with skip connections
\item Benefit from both transfer learning and task-specific loss functions
\item Enable detailed scene understanding
\end{itemize}

\subsection{Performance Comparison}

Table~\ref{tab:overall} provides a cross-laboratory performance summary.

\begin{table}[H]
\centering
\caption{Overall Performance Summary}
\label{tab:overall}
\begin{tabular}{@{}lll@{}}
\toprule
\textbf{Laboratory} & \textbf{Best Method} & \textbf{Best Performance} \\ \midrule
Lab 1 (BoVW) & Dense SIFT + RBF SVM & 81.85\% accuracy \\
Lab 2 (Transfer) & ResNet-18 Fine-tuning & 75.6\% accuracy \\
Lab 3 (Segmentation) & EfficientNet + Dice & 87.1\% mIoU \\ \bottomrule
\end{tabular}
\end{table}

Note: Direct comparison is not meaningful due to different tasks and datasets, but trends are evident.

\subsection{Lessons Learned}

\textbf{1. Feature Representation is Critical:}
\begin{itemize}
\item In Lab 1, switching to Dense SIFT had a larger impact than changing classifiers
\item In Labs 2-3, using pre-trained features dramatically improved results
\end{itemize}

\textbf{2. Transfer Learning is Powerful:}
\begin{itemize}
\item Lab 2 achieved good results on fine-grained classification with only 200 samples per class
\item Lab 3 achieved high-quality segmentation with just 200 training images
\end{itemize}

\textbf{3. Task-Specific Design Matters:}
\begin{itemize}
\item Histogram intersection kernel (Lab 1) is designed for histograms but was outperformed by general-purpose RBF
\item Dice Loss (Lab 3) directly optimized the evaluation metric and improved results
\end{itemize}

\textbf{4. Data Augmentation Improves Generalization:}
\begin{itemize}
\item Both Labs 2 and 3 benefited significantly from augmentation
\item Augmentation is especially important with limited training data
\end{itemize}

\textbf{5. Hyperparameter Tuning is Essential:}
\begin{itemize}
\item Small changes in learning rate, batch size, or architecture can yield noticeable improvements
\item Systematic experimentation (grid search) is valuable
\end{itemize}

\subsection{Practical Implications}

\begin{itemize}
\item \textbf{For new projects:} Start with transfer learning rather than training from scratch
\item \textbf{For limited data:} Use aggressive augmentation and pre-trained models
\item \textbf{For fine-grained tasks:} Fine-tuning outperforms fixed feature extraction
\item \textbf{For dense predictions:} U-Net-style architectures with skip connections work well
\item \textbf{For imbalanced datasets:} Consider task-specific loss functions like Dice Loss
\end{itemize}

\section{Conclusion}

This comprehensive study explored three fundamental aspects of computer vision: classical feature-based image classification, transfer learning for fine-grained recognition, and semantic segmentation with deep learning.

\textbf{Key achievements:}
\begin{itemize}
\item Demonstrated the power of dense feature sampling, achieving 81.85\% accuracy on scene classification with classical methods
\item Successfully applied transfer learning to fine-grained bird classification, achieving 75.6\% accuracy on 200 species
\item Implemented effective semantic segmentation, achieving 87.1\% mIoU on pet segmentation
\end{itemize}

\textbf{Main insights:}
\begin{itemize}
\item Feature representation quality has the largest impact on performance
\item Transfer learning enables strong results with limited labeled data
\item Data augmentation is crucial for generalization
\item Task-specific architectural and loss function choices significantly improve results
\end{itemize}

\textbf{Future directions:}
\begin{itemize}
\item Explore more recent architectures (Vision Transformers, ConvNeXt)
\item Investigate self-supervised pre-training methods
\item Apply these techniques to more challenging real-world datasets
\item Extend to multi-task learning combining classification and segmentation
\end{itemize}

In conclusion, this laboratory series provided hands-on experience with the evolution of computer vision methods, from classical handcrafted features to modern deep learning approaches. The experiments clearly demonstrate why deep learning has become dominant in computer vision: learned representations outperform handcrafted features, transfer learning enables data-efficient learning, and end-to-end optimization produces superior results.

\begin{thebibliography}{00}
\bibitem{b1} J. Sivic and A. Zisserman, ``Video Google: A text retrieval approach to object matching in videos,'' \textit{Proceedings Ninth IEEE International Conference on Computer Vision}, 2003, pp. 1470-1477 vol.2.

\bibitem{b2} D. G. Lowe, ``Distinctive Image Features from Scale-Invariant Keypoints,'' \textit{International Journal of Computer Vision}, vol. 60, no. 2, pp. 91-110, 2004.

\bibitem{b3} K. He, X. Zhang, S. Ren, and J. Sun, ``Deep Residual Learning for Image Recognition,'' \textit{2016 IEEE Conference on Computer Vision and Pattern Recognition (CVPR)}, 2016, pp. 770-778.

\bibitem{b4} O. Ronneberger, P. Fischer, and T. Brox, ``U-Net: Convolutional Networks for Biomedical Image Segmentation,'' \textit{Medical Image Computing and Computer-Assisted Intervention – MICCAI 2015}, pp. 234-241, 2015.

\bibitem{b5} M. Tan and Q. Le, ``EfficientNet: Rethinking Model Scaling for Convolutional Neural Networks,'' \textit{Proceedings of the 36th International Conference on Machine Learning}, 2019.

\bibitem{b6} C. Wah, S. Branson, P. Welinder, P. Perona, and S. Belongie, ``The Caltech-UCSD Birds-200-2011 Dataset,'' Technical Report CNS-TR-2011-001, California Institute of Technology, 2011.

\bibitem{b7} O. M. Parkhi et al., ``Cats and Dogs,'' \textit{2012 IEEE Conference on Computer Vision and Pattern Recognition}, 2012, pp. 3498-3505.

\bibitem{b8} J. Deng, W. Dong, R. Socher, L.-J. Li, K. Li and L. Fei-Fei, ``ImageNet: A Large-Scale Hierarchical Image Database,'' \textit{2009 IEEE Conference on Computer Vision and Pattern Recognition}, 2009, pp. 248-255.

\bibitem{b9} F. Milletari, N. Navab, and S.-A. Ahmadi, ``V-Net: Fully Convolutional Neural Networks for Volumetric Medical Image Segmentation,'' \textit{2016 Fourth International Conference on 3D Vision (3DV)}, 2016, pp. 565-571.

\bibitem{b10} A. Howard et al., ``Searching for MobileNetV3,'' \textit{2019 IEEE/CVF International Conference on Computer Vision (ICCV)}, 2019, pp. 1314-1324.
\end{thebibliography}

\end{document}
